\documentclass[11pt]{article}
\usepackage[a4paper,margin=1in]{geometry}
\usepackage{amsmath,amssymb,amsthm,booktabs}
\usepackage[T1]{fontenc}
\usepackage{lmodern}
\usepackage[hidelinks]{hyperref}
\setlength{\emergencystretch}{3em}

\theoremstyle{plain}
\newtheorem{theorem}{Theorem}
\newtheorem{lemma}[theorem]{Lemma}
\newtheorem{proposition}[theorem]{Proposition}
\newtheorem{corollary}[theorem]{Corollary}
\theoremstyle{definition}
\newtheorem{definition}[theorem]{Definition}
\theoremstyle{remark}
\newtheorem{remark}[theorem]{Remark}

\title{A Proof of Conjecture 00000000049}
\author{%
  Feiyu\thanks{Independent researcher.}
  \and
  Claude (Opus 5)\thanks{Anthropic.}%
}
\date{2026-09-13}

\begin{document}
\maketitle

\begin{abstract}
Conjecture 00000000049 asserts, of point sets in the plane all of whose pairwise
distances are prime, that the maximal size is $4$ and that a set of size $4$
exists. Both hold. A set of size $4$ is $\{0,2,5,7\}$ on a line. For the bound
we show that among the six distances of any such quadruple \emph{exactly two
equal $2$}: the six squared distances satisfy the Gram relation of four coplanar
points, each is $4$ or is $1$ modulo $8$, and a check of the $64$ resulting
patterns leaves only a perfect matching, a triangle, or all six pairs --- the
last two being excluded by the exact value $128$ and by a computation modulo
$32$. Five points would then give $3e=5\cdot2$ for the number $e$ of pairs at
distance $2$, which is impossible. The whole conjecture, both parts, is
formalized in Lean~4 against Mathlib.
\end{abstract}

\section{Statement and conventions}

We identify the Euclidean plane with $\mathbb C$. Call a finite set
$S\subset\mathbb C$ a \emph{prime-distance set} if $|P-Q|$ is a prime integer
for all distinct $P,Q\in S$. Conjecture 00000000049 asserts
\begin{enumerate}
  \item[(a)] every prime-distance set has at most $4$ elements;
  \item[(b)] some prime-distance set has $4$ elements.
\end{enumerate}

Throughout, ``prime'' means a rational prime $2,3,5,7,\ldots$; note that $2$ is
allowed, and it is the only even one. That is what makes the problem differ from
the classical question about odd integral distances, and, as it turns out, the
distance $2$ is what the whole proof is about.

\section{Part (b)}

\begin{proposition}\label{prop:b}
$S=\{0,2,5,7\}\subset\mathbb R\subset\mathbb C$ is a prime-distance set of size
$4$.
\end{proposition}

\begin{proof}
The six distances are
\[
  |0-2|=2,\quad |0-5|=5,\quad |0-7|=7,\quad
  |2-5|=3,\quad |2-7|=5,\quad |5-7|=2,
\]
all prime. \end{proof}

\section{Part (a)}

\subsection{The Gram relation}

Let $P_0,P_1,P_2,P_3$ be four points of the plane, write $v_i=P_i-P_0$, and put
\[
  s_i=|P_0P_i|^2,\qquad t_{ij}=|P_iP_j|^2 .
\]
The Gram matrix $G=(v_i\cdot v_j)_{1\le i,j\le3}$ is singular, because
$v_1,v_2,v_3$ lie in a $2$-dimensional space. Since
$2\,v_i\cdot v_j=s_i+s_j-t_{ij}$ for $i\ne j$ and $2\,v_i\cdot v_i=2s_i$,
setting
\[
  x_{ij}=s_i+s_j-t_{ij}
\]
and expanding $\det(2G)=0$ gives, after dividing by $2$,
\begin{equation}\label{eq:E}
  E:=4s_1s_2s_3+x_{12}x_{13}x_{23}
     -s_1x_{23}^2-s_2x_{13}^2-s_3x_{12}^2=0 .
\end{equation}
As a polynomial identity in the coordinates of the four points this is the
statement that the Gram determinant of three coplanar vectors vanishes; it is
$\pm\frac12$ times the Cayley--Menger determinant of the quadruple
\cite{Blumenthal}, and in particular it is symmetric under relabelling the
points.

Everything below is an arithmetic consequence of \eqref{eq:E}. Note that the six
squared distances of a prime-distance quadruple are integers, and
\begin{equation}\label{eq:mod8}
  p^2\equiv
  \begin{cases}
    4 \pmod 8, & p=2,\\
    1 \pmod 8, & p\ \text{an odd prime},
  \end{cases}
\end{equation}
since an odd square is $\equiv1\pmod 8$.

\subsection{Which pairs can be at distance two}

Let $S$ be a prime-distance set of size $4$ and let $H$ be the graph on $S$
whose edges are the pairs at distance exactly $2$. By \eqref{eq:mod8} the
residue of $E$ modulo $8$ depends only on $H$, so it can be computed once for
each of the eleven graphs on four vertices.

\begin{lemma}\label{lem:mod8}
$H$ is a perfect matching (two disjoint edges), a triangle together with an
isolated vertex, or the complete graph.
\end{lemma}

\begin{proof}
Substituting $4$ for the squared distances of the edges of $H$ and $1$ for the
others in \eqref{eq:E} and reducing modulo $8$ gives, for the eleven isomorphism
types on four vertices:

\begin{center}
\begin{tabular}{lcc}
\toprule
$H$ & edges & $E \bmod 8$\\
\midrule
empty                       & 0 & 2\\
one edge                    & 1 & 4\\
two adjacent edges          & 2 & 7\\
\textbf{perfect matching}   & 2 & \textbf{0}\\
path $P_4$                  & 3 & 7\\
star $K_{1,3}$              & 3 & 3\\
\textbf{triangle $+$ isolated vertex} & 3 & \textbf{0}\\
$4$-cycle                   & 4 & 6\\
triangle $+$ pendant edge   & 4 & 4\\
$K_4$ minus an edge         & 5 & 4\\
\textbf{$K_4$}              & 6 & \textbf{0}\\
\bottomrule
\end{tabular}
\end{center}

All but the three boldface types have $E\not\equiv0\pmod 8$, contradicting
\eqref{eq:E}. (The computation is a finite check over the $2^6=64$ labelled
patterns; it is the Lean lemma \texttt{pattern\_mod8}.) \end{proof}

\begin{lemma}\label{lem:k4}
$H$ is not the complete graph.
\end{lemma}

\begin{proof}
All six squared distances would equal $4$, and then
$E=4\cdot64+64-3\cdot4\cdot16=128\ne0$. \end{proof}

\begin{lemma}\label{lem:tri}
$H$ is not a triangle together with an isolated vertex.
\end{lemma}

\begin{proof}
Take the isolated vertex as the base point $P_0$, so that $t_{12}=t_{13}=t_{23}=4$
while $s_1=p^2$, $s_2=q^2$, $s_3=r^2$ for odd primes $p,q,r$. The square of an
odd number is $1$, $9$, $17$ or $25$ modulo $32$, and for each of the $4^3=64$
choices \eqref{eq:E} gives
\[
  E\equiv16\pmod{32},
\]
never $0$. (Again a finite check: the Lean lemma \texttt{tri\_mod32}, which
covers the four positions of the isolated vertex.) \end{proof}

\begin{theorem}\label{thm:two}
In a prime-distance set of size $4$, exactly two of the six distances equal $2$,
and those two pairs are disjoint.
\end{theorem}

\begin{proof}
Lemmas~\ref{lem:mod8}, \ref{lem:k4} and~\ref{lem:tri}. \end{proof}

\subsection{No five points}

\begin{theorem}\label{thm:five}
Every prime-distance set has at most $4$ elements.
\end{theorem}

\begin{proof}
Suppose $S$ is a prime-distance set with $|S|\ge5$ and pick five points
$Q_1,\ldots,Q_5$ of it. Let $e$ be the number of pairs among them at distance
$2$. Each of the five $4$-element subsets is a prime-distance set, so by
Theorem~\ref{thm:two} it contains exactly two such pairs; and each pair lies in
exactly $\binom32=3$ of those subsets. Counting the incidences between pairs at
distance $2$ and $4$-element subsets in two ways,
\[
  3e=5\cdot2=10 ,
\]
which has no solution in integers. \end{proof}

\begin{corollary}
Conjecture 00000000049 holds.
\end{corollary}

\begin{proof}
Theorem~\ref{thm:five} and Proposition~\ref{prop:b}. \end{proof}

\begin{remark}
The bound needs no geometry beyond \eqref{eq:E}: no case analysis on
configurations, no coordinates, and --- in contrast with the usual route through
the Erd\H os--Anning theorem \cite{Anning} --- no collinearity. Collinearity is nevertheless
true here, and Appendix~\ref{sec:appendix} proves it, together with the fact
that $\{0,2,5,7\}$ is the \emph{only} quadruple up to isometry.
\end{remark}

\section{Formal verification}

The accompanying Lean~4 project proves the conjecture in full: the final theorem
is
\[
  \texttt{conjecture\_holds} : \texttt{ConjectureHolds},
\]
where \texttt{ConjectureHolds} is the conjecture itself, stated in the
formalization as
\[
  \bigl(\forall n\ (P:\texttt{Fin}\ n\to\mathbb C),\
    \texttt{IsPrimeDistanceSet}\ P\to n\le4\bigr)
  \ \wedge\
  \bigl(\exists P:\texttt{Fin}\ 4\to\mathbb C,\ \texttt{IsPrimeDistanceSet}\ P\bigr).
\]

\begin{itemize}
  \item \texttt{IsPrimeDistanceSet}, \texttt{ConjectureHolds} --- the notion and
    the conjecture, with the plane modeled as $\mathbb C$ and primality
    Mathlib's \texttt{Nat.Prime};
  \item \texttt{witness}, \texttt{dTable\_prime}, \texttt{conjecture\_part\_b}
    --- Proposition~\ref{prop:b};
  \item \texttt{gram}, \texttt{gram\_eq\_zero}, \texttt{gram\_dist},
    \texttt{gram\_int} --- the relation \eqref{eq:E}. The identity itself is
    proved by \texttt{ring} on the coordinates of the three difference vectors;
  \item \texttt{pattern\_mod8} --- Lemma~\ref{lem:mod8}, as a \texttt{decide}
    over the $64$ labelled patterns;
  \item \texttt{tri\_mod32}, \texttt{odd\_sq\_mod32} ---
    Lemma~\ref{lem:tri}, likewise by \texttt{decide};
  \item \texttt{two\_of\_six}, \texttt{two\_of\_four} ---
    Theorem~\ref{thm:two}, in the counting form used next: the six indicators
    of ``distance $=2$'' sum to $2$;
  \item \texttt{no\_five}, \texttt{card\_le\_four} --- Theorem~\ref{thm:five};
    the double count is discharged by \texttt{omega} from the five equations;
  \item \texttt{conjecture\_holds} --- the conjecture;
  \item \texttt{gaps\_of\_four\_collinear} --- the arithmetic half of
    Appendix~\ref{sec:appendix} (Lemma~\ref{lem:gaps}).
\end{itemize}

\paragraph{Scope.} Parts (a) and (b) are both formalized, with no gap: the Lean
development contains no \texttt{sorry}, no \texttt{admit} and no
\texttt{native\_decide}, and \texttt{\#print axioms conjecture\_holds} reports
exactly the three standard axioms \texttt{propext}, \texttt{Classical.choice}
and \texttt{Quot.sound}. What is \emph{not} formalized is the supplementary
Appendix~\ref{sec:appendix}: its geometric step (Lemma~\ref{lem:coll}) is a
coordinate computation that would need a rigid-motion normalization in Lean, and
nothing in the conjecture depends on it. The project builds with \texttt{lake
build} on \texttt{leanprover/lean4:v4.33.1} with Mathlib \texttt{v4.33.1}.

\section{Computational corroboration}

Four points are planar exactly when their $5\times5$ Cayley--Menger determinant
vanishes; with prime squared distances this is an exact integer test. Over all
$6$-tuples of primes below $120$ satisfying the triangle inequalities --- about
$4.4\cdot10^6$ configurations --- the only solutions are degenerate, and every
one of them is $\{0,2,5,7\}$ up to relabelling and reflection. This agrees with
Theorem~\ref{thm:two} and with Appendix~\ref{sec:appendix}. The search is
corroboration, not part of the proof.

\appendix

\section{The quadruples are collinear, and there is only one}
\label{sec:appendix}

Theorem~\ref{thm:two} says that a prime-distance quadruple has exactly two
distances equal to $2$, on disjoint pairs. That is already enough for the
conjecture. With one coordinate computation more, the quadruples can be
determined completely. Nothing in this appendix is used above, and it is not
formalized apart from Lemma~\ref{lem:gaps}.

\begin{lemma}\label{lem:coll}
Let $\{A,B,C,D\}$ be a prime-distance set of size $4$ with $|AB|=|CD|=2$ and the
other four distances odd. Then the four points are collinear.
\end{lemma}

\begin{proof}
Place $A=(-1,0)$, $B=(1,0)$. For any point $P=(x,y)$ with $|AP|,|BP|$ odd,
\[
  x=\frac{|AP|^2-|BP|^2}{4}
\]
is an even integer, because $|AP|^2-|BP|^2\equiv0\pmod 8$; and
$y^2=|AP|^2-(x+1)^2\equiv1-1=0\pmod 8$. Apply this to $C=(x_C,y_C)$ and
$D=(x_D,y_D)$.

From $|CD|^2=4$ and $x_C-x_D$ even we get $(x_C-x_D)^2\in\{0,4\}$.

If $x_C=x_D$ then $(y_C-y_D)^2=4$, so $2y_Cy_D=y_C^2+y_D^2-4\equiv4\pmod8$; thus
$y_Cy_D$ is an integer $\equiv2\pmod4$ and $(y_Cy_D)^2\equiv4\pmod{32}$. But
$(y_Cy_D)^2=y_C^2y_D^2\equiv0\pmod{64}$. Contradiction.

So $|x_C-x_D|=2$ and $y_C=y_D=:y$. Relabelling, write $C=(x,y)$, $D=(x+2,y)$ and
put $p=|AC|$, $q=|AD|$, $r=|BC|$; then $|BD|=|AC|=p$ and
\[
  q^2+r^2=2p^2+8,\qquad x=\frac{p^2-r^2}{4},\qquad y^2=p^2-(x+1)^2 .
\]
Now $y^2\ge0$ forces $(p-2)^2\le r^2\le(p+2)^2$, i.e. $p-2\le r\le p+2$; as $p,r$
are odd, $r\in\{p-2,p,p+2\}$. If $r=p$ then $q^2=p^2+8$, so $(q-p)(q+p)=8$ with
both factors even, forcing $q-p=2$, $q+p=4$ and $p=1$, not prime. Hence
$r=p\pm2$, which gives $x+1=\pm p$ and $y^2=0$: the four points lie on the
$x$-axis. \end{proof}

\begin{lemma}\label{lem:gaps}
Let $x_0<x_1<x_2<x_3$ be four points on a line with all six pairwise distances
prime, and let $g_i=x_i-x_{i-1}$. Then $(g_1,g_2,g_3)=(2,3,2)$.
\end{lemma}

\begin{proof}
Each $g_i$ is prime, and so is each sum of consecutive ones. Two adjacent gaps
cannot both be odd, since their sum would be even and at least $6$; and they
cannot both be $2$, since $4$ is not prime. So of each adjacent pair exactly one
equals $2$, which leaves the patterns $(2,\text{odd},2)$ and
$(\text{odd},2,\text{odd})$. In the second the total span $g_1+g_2+g_3$ is even
and at least $8$, hence not prime; so $g_1=g_3=2$ and $q:=g_2$ is odd.

Now $q$, $q+2$ and $q+4$ are all prime --- they are $|x_1x_2|$, $|x_1x_3|$ and
$|x_0x_3|$ --- and one of any three integers in arithmetic
progression of difference $2$ is divisible by $3$. A prime divisible by $3$
equals $3$; since $q\ge3$ the values $q+2$ and $q+4$ exceed $3$, so the one
divisible by $3$ is $q$ itself, and $q=3$. \end{proof}

\begin{theorem}
Every prime-distance set of size $4$ is collinear, and is $\{0,2,5,7\}$ up to an
isometry of the plane.
\end{theorem}

\begin{proof}
By Theorem~\ref{thm:two} the two distances equal to $2$ lie on disjoint pairs, so
Lemma~\ref{lem:coll} applies and the set is collinear; by Lemma~\ref{lem:gaps}
its consecutive gaps are $2,3,2$, so after an isometry of the line carrying it
into $\mathbb R$ it is $\{0,2,5,7\}$. \end{proof}

\begin{remark}
The same gap analysis reproves, in one line, that no five points on a line have
prime pairwise distances: the four gaps would have to alternate between $2$ and
an odd prime, making the span $4$ plus two odd primes --- even and at least
$10$.
\end{remark}

\begin{thebibliography}{9}

\bibitem{Anning} N.~H.~Anning, P.~Erd\H os,
\emph{Integral distances},
Bull. Amer. Math. Soc. \textbf{51} (1945), 598--600.

\bibitem{Blumenthal} L.~M.~Blumenthal,
\emph{Theory and Applications of Distance Geometry},
Clarendon Press, Oxford, 1953.

\end{thebibliography}

\end{document}
